%%%%%%%%%%%%%%%%%%%%%%%%%%%%%%%%%
% Results and Findings          %
%%%%%%%%%%%%%%%%%%%%%%%%%%%%%%%%%

\section{RESULTS AND FINDINGS}
\label{sec:results}

\noindent
This section presents the results and findings of the study. It first consolidates the experimental framework that the four iterations of Section~\ref{sec:iterations} produced, defines the evaluation metrics and protocols, and then reports the measurements as direct answers to the three research questions posed in Section~\ref{sec:introduction}. The interpretation of these results, including the mechanisms behind them, the trade-offs involved, and the study's limitations, is deferred to the Discussion in Section~\ref{sec:discussion}.

% ═══════════════════════════════════════════════════════════
\subsection{Presentation of the Experimental Framework}
\label{subsec:framework}

\noindent
The framework evaluated in this section is a supervised inverse-synthesis system that maps a rendered audio clip to the 16 Vital parameters that produced it. It was built and refined across four iterations, and its components are summarised below with a pointer to the iteration in which each was designed and implemented.

\begin{itemize}
    \item[$\blacksquare$] \textbf{Data generation pipeline.} The timbre-constrained preset generation, audio rendering, silence filtering, multi-scale mel extraction, and preset-level splitting that produce the training corpus were built in Iteration~1 (Figure~\ref{fig:pipeline}).
    \item[$\blacksquare$] \textbf{Model family.} The three backbone architectures compared throughout, the residual CNN (\texttt{SynthParamNet}), AST (ViT), and ConvNeXt, together with the sinusoidal pitch-embedding branch, were established in Iteration~2 (Figure~\ref{fig:v1_design}).
    \item[$\blacksquare$] \textbf{Architectural enhancements.} The three-channel multi-scale mel input, FiLM pitch conditioning, the grouped prediction head, and the group-weighted SmoothL1 objective were introduced in Iteration~3 (Figure~\ref{fig:v2_design}).
    \item[$\blacksquare$] \textbf{Data regime.} The multi-pitch re-rendering of every preset and the online SpecAugment regularisation that define the enlarged corpus were introduced in Iteration~4 (Figure~\ref{fig:v3_design}).
\end{itemize}

\noindent
The three iterations that produced trained models are referred to as V1 (single-scale baseline), V2 (enhanced architecture on the single-pitch corpus), and V3 (enhanced architecture on the multi-pitch corpus). These versions are the experimental conditions from which the results below are drawn; each research question is answered using the version comparisons relevant to it.

\subsubsection{Evaluation Metrics}

\noindent
Model quality is assessed with the three metrics defined in Table~\ref{tab:metrics}. The normalised MAE is the primary accuracy metric; the per-parameter MAE localises that accuracy to individual synthesis parameters; and the Multi-Scale Spectral (MSS) loss provides an audio-domain check by re-rendering predictions through Vital.

\begin{table}[H]
\centering
\renewcommand{\arraystretch}{1.4}
\begin{tabular}{p{0.20\linewidth} p{0.42\linewidth} p{0.28\linewidth}}
\hline
\textbf{Metric} & \textbf{Description} & \textbf{Significance} \\
\hline
Normalised MAE &
Mean absolute error between predicted and ground-truth parameter vectors in the normalised $[0,1]$ space, averaged over all 16 parameters. &
Primary accuracy metric; lower is better. \\
Per-parameter MAE &
The same error computed separately for each of the 16 parameters. &
Reveals which synthesis parameters are hard to recover from audio. \\
Multi-Scale Spectral (MSS) loss &
Mean absolute log-mel-spectrogram difference across four FFT scales between the original clip and the audio re-rendered from the predicted parameters. &
Audio-domain check that parameter accuracy translates to perceptual fidelity; lower is better. \\
\hline
\end{tabular}
\vspace{0.3cm}
\caption{Evaluation metrics used in this study. MAE is a parameter-space metric; MSS is an audio-space metric computed by re-rendering predictions through Vital (evaluation only, since Vital is non-differentiable).}
\label{tab:metrics}
\end{table}

\subsubsection{Evaluation Protocols}

\noindent
Two evaluation protocols are used and kept distinct throughout. V1 and V2 are each reported on their own five-fold cross-validation, and V3 on its own single preset-level 80/10/10 split, so that every version is first interpreted on the protocol under which it was trained. In addition, because the MSS metric requires re-rendering every prediction through Vital, all \emph{direct} V2-versus-V3 comparisons of MAE and MSS are computed on a fixed shared set of 100 held-out clips evaluated identically for both versions. Each table and figure states which protocol it uses.

% ═══════════════════════════════════════════════════════════
\subsection{Answers to the Research Questions}
\label{subsec:rq_answers}

% -----------------------------------------------------------
\subsubsection{Research Question 1}
\label{subsec:rq1}

\noindent
\textbf{Research Question 1:} \textit{How do convolutional architectures (residual CNN and ConvNeXt) compare to AST in regressing synthesizer parameters from mel-spectrograms?}

\noindent \newline
Each of the three backbone families (SynthParamNet, AST-small, and ConvNeXt-tiny) was trained under three pitch conditioning strategies (no pitch, basic MLP, and sinusoidal embedding), yielding nine V1 configurations evaluated over five stratified folds. Table~\ref{tab:v1_results} reports the five-fold test MAE for all nine configurations, with the AST-small no-pitch variant as the reference point that approximates the Bruford et al.\ setup \cite{bruford2024ast}.

\begin{table}[H]
\centering
\renewcommand{\arraystretch}{1.3}
\begin{tabular}{llcc}
\hline
\textbf{Backbone} & \textbf{Pitch Conditioning} & \textbf{MAE (mean $\pm$ std)} & \textbf{$\Delta$ vs baseline} \\
\hline
\multirow{3}{*}{SynthParamNet (CNN)}
  & No pitch      & $0.0850 \pm 0.0014$ & $+0.0057$ \\
  & Basic MLP     & $0.0828 \pm 0.0007$ & $+0.0035$ \\
  & Sinusoidal    & $0.0829 \pm 0.0008$ & $+0.0036$ \\
\hline
\multirow{3}{*}{AST-small}
  & No pitch$^{\dagger}$ & $0.0793 \pm 0.0006$ & --- \\
  & Basic MLP     & $0.0781 \pm 0.0005$ & $-0.0012$ \\
  & Sinusoidal    & $0.0776 \pm 0.0007$ & $-0.0017$ \\
\hline
\multirow{3}{*}{ConvNeXt-tiny}
  & No pitch      & $0.0767 \pm 0.0004$ & $-0.0026$ \\
  & Basic MLP     & $0.0765 \pm 0.0007$ & $-0.0028$ \\
  & \textbf{Sinusoidal} & $\mathbf{0.0758 \pm 0.0005}$ & $\mathbf{-0.0035}$ \\
\hline
\end{tabular}
\vspace{0.3cm}
\caption{V1 five-fold test MAE across all nine backbone $\times$ pitch conditioning configurations. $\Delta$ is computed relative to the AST-small no-pitch baseline ($^{\dagger}$), which approximates the Bruford et al.\ \cite{bruford2024ast} setup. Best result in bold.}
\label{tab:v1_results}
\end{table}

\noindent
Table~\ref{tab:v1_results} and Figure~\ref{fig:v1_architecture} show three results. ConvNeXt-tiny records a lower MAE than AST-small in every pitch condition, by $0.0016$ to $0.0026$. The SynthParamNet CNN is the weakest backbone in all three conditions. Adding pitch conditioning lowers the MAE of every backbone relative to its no-pitch variant, and the sinusoidal embedding gives the lowest error for AST and ConvNeXt, trailing the basic MLP only on the CNN and there by $0.0001$. The best V1 configuration is ConvNeXt-tiny with the sinusoidal embedding at $0.0758 \pm 0.0005$.

\begin{figure}[H]
\centering
\includegraphics[width=\linewidth]{figs/fig_v1_architecture.pdf}
\caption{V1 five-fold test MAE for all nine configurations. The dotted red line marks the AST-small no-pitch reference. ConvNeXt-tiny with sinusoidal embedding achieves the lowest MAE in V1.}
\label{fig:v1_architecture}
\end{figure}

\noindent
The V1 sweep establishes the ranking at the smaller corpus, but the architecture question spans all three versions and both metrics. Table~\ref{tab:arch_compare} reports the best ConvNeXt-tiny and AST-small model at each version on parameter-space MAE and, where it can be measured, on the audio-space MSS loss.

\begin{table}[H]
\centering
\renewcommand{\arraystretch}{1.3}
\begin{tabular}{llcccc}
\hline
\textbf{Version} & \textbf{Corpus} & \textbf{ConvNeXt MAE} & \textbf{AST MAE} & \textbf{ConvNeXt MSS} & \textbf{AST MSS} \\
\hline
V1 & 42k  & $\mathbf{0.0758}$ & $0.0776$ & --- & --- \\
V2 & 42k  & $\mathbf{0.0739}$ & $0.0755$ & $1.405$ & $\mathbf{1.292}$ \\
V3 & 168k & $0.0530$ & $\mathbf{0.0525}$ & $1.222$ & $\mathbf{1.096}$ \\
\hline
\end{tabular}
\vspace{0.3cm}
\caption{ConvNeXt-tiny versus AST-small at each version, using the best pitch-conditioned configuration of each (V1: sinusoidal; V2, V3: FiLM). V1 MAE is the five-fold test average; V2 and V3 MAE and MSS are on the shared 100-sample set. The lower (better) value in each pair is in bold. MSS requires re-rendering through Vital and is unavailable for V1.}
\label{tab:arch_compare}
\end{table}

\noindent
The comparison is both scale-dependent and metric-dependent. On parameter-space MAE, ConvNeXt-tiny is the lower of the two in V1 and V2, and AST-small is lower in V3, by a narrow $0.0005$. On the audio-space MSS metric, AST-small is the lower in both V2 and V3, including V2, where its MAE is higher than ConvNeXt's. ConvNeXt-tiny therefore leads on parameter error at the smaller corpus, whereas AST-small leads on perceptual reconstruction whenever MSS is measured and overtakes ConvNeXt on both metrics at the larger V3 corpus. The best V3 model, AST-small with FiLM, reaches $0.0525$ MAE and $1.096$ MSS.

% -----------------------------------------------------------
\subsubsection{Research Question 2}
\label{subsec:rq2}

\noindent
\textbf{Research Question 2:} \textit{Does replacing a single mel-spectrogram with a three-scale representation combining fine, mid, and coarse spectrograms improve synthesizer parameter estimation accuracy?}

\noindent \newline
The V2 architecture introduces four joint enhancements over V1: a three-channel multi-scale mel input (fine, mid, and coarse), FiLM-based pitch conditioning, a grouped prediction head, and a group-weighted SmoothL1 training objective. To separate the contribution of pitch conditioning from the remaining enhancements, each V2 backbone is trained in both a full FiLM variant and a no-pitch variant that removes the FiLM branch while retaining all other V2 components. Table~\ref{tab:v2_results} reports the five-fold test MAE for the four V2 configurations alongside the V1 best and baseline.

\begin{table}[H]
\centering
\renewcommand{\arraystretch}{1.3}
\begin{tabular}{llcc}
\hline
\textbf{Version} & \textbf{Model} & \textbf{MAE (mean $\pm$ std)} & \textbf{$\Delta$ vs V1 baseline} \\
\hline
V1 & AST-small, no pitch (baseline) & $0.0793 \pm 0.0006$ & --- \\
V1 & ConvNeXt-tiny, sinusoidal (V1 best) & $0.0758 \pm 0.0005$ & $-4.4\%$ \\
\hline
V2 & ConvNeXt-tiny, no pitch & $0.0681 \pm 0.0002$ & $-14.1\%$ \\
V2 & AST-small, no pitch & $0.0727 \pm 0.0003$ & $-8.3\%$ \\
V2 & AST-small, FiLM & $0.0713 \pm 0.0003$ & $-10.1\%$ \\
V2 & \textbf{ConvNeXt-tiny, FiLM} & $\mathbf{0.0665 \pm 0.0002}$ & $\mathbf{-16.1\%}$ \\
\hline
\end{tabular}
\vspace{0.3cm}
\caption{Five-fold test MAE for V1 reference configurations and all four V2 variants. The ConvNeXt-tiny no-pitch row isolates the contribution of the non-pitch V2 components (multi-scale input, grouped head, weighted loss) from the FiLM conditioning effect.}
\label{tab:v2_results}
\end{table}

\noindent
The no-pitch V2 ConvNeXt model ($0.0681 \pm 0.0002$) improves on the best V1 model ($0.0758 \pm 0.0005$) by $10.2\%$ without any pitch input, isolating the contribution of the three non-pitch V2 components (multi-scale input, grouped head, weighted loss). Adding FiLM lowers MAE by a further $0.0016$ ($2.4\%$) for ConvNeXt and $0.0014$ ($1.9\%$) for AST, both consistent across all five folds, giving the best V2 model at $0.0665 \pm 0.0002$ ($-16.1\%$ versus the V1 baseline).

\noindent \newline
Figure~\ref{fig:convergence} plots the training SmoothL1 loss and the validation MAE per epoch (two different quantities). For V2, the training loss falls below roughly $0.015$ by epoch 40 while the validation MAE stops improving and plateaus near $0.066$ for ConvNeXt and near $0.070$ for AST; the V3 panel shows this gap between the two curves substantially narrowed.

\begin{figure}[H]
\centering
\includegraphics[width=\linewidth]{figs/fig_convergence.pdf}
\caption{Training convergence for V2 (ConvNeXt-tiny + FiLM, split~0) and V3 (AST-small + FiLM), each showing the per-epoch training SmoothL1 loss and the validation MAE (two different quantities). The shaded region on the V2 panel highlights the gap between the falling training loss and the plateaued validation MAE that motivates the dataset expansion in V3.}
\label{fig:convergence}
\end{figure}

% -----------------------------------------------------------
\subsubsection{Research Question 3}
\label{subsec:rq3}

\noindent
\textbf{Research Question 3:} \textit{Under what conditions does FiLM-based pitch conditioning improve synthesizer parameter estimation, and why does its effectiveness depend on the training data?}

\noindent \newline
The FiLM-conditioned versus no-pitch ablation from V2 is repeated under the V3 training regime: the same $\approx 42{,}250$ presets re-rendered at up to three fixed MIDI pitches each ($\approx 168{,}000$ samples total after silence filtering), with SpecAugment regularisation applied online. All four V3 models are trained on a single 80/10/10 preset-level split, so V3 MAE values are reported without cross-validation standard deviations.

\paragraph{FiLM conditioning gap: V2 vs V3.}
Table~\ref{tab:film_gap} reports the FiLM conditioning gap for both backbones under both data regimes, and Figure~\ref{fig:film_gap} visualises it. The V3 rows use the shared 100-sample set so that the V3 gap is measured on the same protocol as the cross-version comparison in Table~\ref{tab:arch_compare}.

\begin{table}[H]
\centering
\renewcommand{\arraystretch}{1.3}
\begin{tabular}{llllcc}
\hline
\textbf{Version} & \textbf{Backbone} & \textbf{Dataset} & \textbf{Pitch} & \textbf{MAE} & \textbf{FiLM gain} \\
\hline
\multirow{4}{*}{V2}
  & \multirow{2}{*}{ConvNeXt-tiny} & \multirow{2}{*}{42k, 5-fold} & FiLM     & $0.0665 \pm 0.0002$ & \multirow{2}{*}{$-0.0016$ ($-2.4\%$)} \\
  &                                &                              & No pitch & $0.0681 \pm 0.0002$ & \\
  \cline{2-6}
  & \multirow{2}{*}{AST-small}     & \multirow{2}{*}{42k, 5-fold} & FiLM     & $0.0713 \pm 0.0003$ & \multirow{2}{*}{$-0.0014$ ($-1.9\%$)} \\
  &                                &                              & No pitch & $0.0727 \pm 0.0003$ & \\
\hline
\multirow{4}{*}{V3}
  & \multirow{2}{*}{ConvNeXt-tiny} & \multirow{2}{*}{168k, shared} & FiLM    & $0.0530$ & \multirow{2}{*}{$-0.0020$ ($-3.6\%$)} \\
  &                                &                               & No pitch & $0.0550$ & \\
  \cline{2-6}
  & \multirow{2}{*}{AST-small}     & \multirow{2}{*}{168k, shared} & \textbf{FiLM}    & $\mathbf{0.0525}$ & \multirow{2}{*}{$-0.0035$ ($-6.3\%$)} \\
  &                                &                               & No pitch & $0.0560$ & \\
\hline
\end{tabular}
\vspace{0.3cm}
\caption{FiLM conditioning gain for ConvNeXt-tiny and AST-small under V2 (single-pitch, five-fold CV) and V3 (multi-pitch, shared 100-sample set). The gain grows from V2 to V3 for both backbones. On the strictly matched shared set the V2 ConvNeXt gain does not reproduce (see Table~\ref{tab:arch_compare}); the reliability of the two protocols is discussed in Section~\ref{sec:discussion}.}
\label{tab:film_gap}
\end{table}

\begin{figure}[H]
\centering
\includegraphics[width=\linewidth]{figs/fig_film_gap.pdf}
\caption{FiLM vs no-pitch MAE for both backbones in V2 (five-fold) and V3 (shared set). The annotated $\Delta$ values show the FiLM gain growing from V2 to V3, particularly for AST-small.}
\label{fig:film_gap}
\end{figure}

\noindent
The FiLM gain grows from V2 to V3 for both backbones: from $-0.0016$ ($2.4\%$) to $-0.0020$ ($3.6\%$) for ConvNeXt, and from $-0.0014$ ($1.9\%$) to $-0.0035$ ($6.3\%$) for AST. On the shared set, V3 FiLM improves both backbones on both MAE and MSS (Table~\ref{tab:arch_compare}). Because V3 uses a single split, these gains carry no cross-validation variance, and the two protocols disagree on the sign of the V2 ConvNeXt gain; both points are taken up in Section~\ref{sec:discussion}.

\paragraph{Cross-iteration summary.}
Reading Table~\ref{tab:arch_compare} across versions, the best V3 model reduces MAE by $29.0\%$ over the best V2 model on the shared set ($0.0739 \rightarrow 0.0525$) and lowers MSS in parallel ($1.405 \rightarrow 1.096$). Relative to the V1 baseline ($0.0793$), the best V3 model, AST-small with FiLM, is a $33.8\%$ MAE reduction. The corresponding shift in the ConvNeXt-versus-AST ranking is reported under RQ1 (Table~\ref{tab:arch_compare}).

\paragraph{Perceptual evaluation across FFT scales.}
Figure~\ref{fig:mss_scales} breaks down the MSS loss across the four evaluation FFT scales for all four V2 and V3 models. Each V3 model attains a lower MSS than its V2 counterpart at every scale, from the energy-sensitive $n_{\text{fft}}=128$ band to the harmonic-sensitive $n_{\text{fft}}=8{,}192$ band.

\begin{figure}[H]
\centering
\includegraphics[width=\linewidth]{figs/fig_mss_scales.pdf}
\caption{Multi-Scale Spectral loss across four FFT scales for all V2 and V3 models (shared eval set, $N=100$). Solid lines: V3 models. Dashed lines: V2 models. Lower is better.}
\label{fig:mss_scales}
\end{figure}

\noindent
One observation is recorded here for the discussion that follows: in V2 the no-pitch variants attain a marginally lower MSS than their FiLM counterparts (ConvNeXt no-pitch $1.383$ vs FiLM $1.405$; AST no-pitch $1.273$ vs FiLM $1.292$), even though FiLM is superior on parameter-space MAE. This reversal does not occur in V3, where the FiLM models attain the lower MSS as well.

\paragraph{Per-parameter analysis.}
Table~\ref{tab:per_param} and Figure~\ref{fig:per_param} report the per-parameter MAE for the best model at each iteration: V1 (ConvNeXt-tiny, sinusoidal), V2 (ConvNeXt-tiny, FiLM), and V3 (AST-small, FiLM). The V1 and V2 columns are the per-parameter breakdown of each version's best cross-validation split, while the V3 column is measured on the shared 100-sample set; the columns are therefore read for the relative difficulty of parameters rather than as an exact like-for-like delta.

\begin{table}[H]
\centering
\renewcommand{\arraystretch}{1.25}
\begin{tabular}{llccc}
\hline
\textbf{Group} & \textbf{Parameter} & \textbf{V1 best} & \textbf{V2 best} & \textbf{V3 best} \\
\hline
\multirow{7}{*}{Oscillator}
  & \texttt{osc1\_wave\_frame}    & 0.1362 & 0.1209 & 0.1015 \\
  & \texttt{osc1\_level}          & 0.1161 & 0.0989 & 0.0823 \\
  & \texttt{osc1\_unison\_voices} & 0.1067 & 0.1010 & 0.1022 \\
  & \texttt{osc1\_unison\_detune} & 0.0869 & 0.0789 & 0.0681 \\
  & \texttt{osc2\_wave\_frame}    & 0.1610 & 0.1508 & 0.1205 \\
  & \texttt{osc2\_level}          & 0.1288 & 0.1081 & 0.0868 \\
  & \texttt{osc2\_transpose}      & 0.0902 & 0.0622 & 0.0308 \\
\hline
\multirow{4}{*}{Envelope}
  & \texttt{env1\_attack}         & 0.0245 & 0.0184 & 0.0117 \\
  & \texttt{env1\_decay}          & 0.0191 & 0.0155 & 0.0103 \\
  & \texttt{env2\_attack}         & 0.0367 & 0.0343 & 0.0237 \\
  & \texttt{env2\_decay}          & 0.0269 & 0.0277 & 0.0213 \\
\hline
\multirow{2}{*}{Filter}
  & \texttt{filter1\_cutoff}      & 0.0415 & 0.0378 & 0.0294 \\
  & \texttt{filter1\_resonance}   & 0.0679 & 0.0665 & 0.0498 \\
\hline
\multirow{3}{*}{Misc}
  & \texttt{sample\_level}        & 0.0809 & 0.0583 & 0.0457 \\
  & \texttt{mod1\_amount}         & 0.0346 & 0.0359 & 0.0290 \\
  & \texttt{reverb\_mix}          & 0.0461 & 0.0426 & 0.0266 \\
\hline
\end{tabular}
\vspace{0.3cm}
\caption{Per-parameter MAE for the best model at each iteration. V1 best: ConvNeXt-tiny with sinusoidal embedding. V2 best: ConvNeXt-tiny with FiLM. V3 best: AST-small with FiLM. V1 and V2 are the per-parameter breakdown of each version's best five-fold split; V3 is measured on the shared 100-sample set.}
\label{tab:per_param}
\end{table}

\begin{figure}[H]
\centering
\includegraphics[width=\linewidth]{figs/fig_per_param.pdf}
\caption{Per-parameter MAE across iterations. Background shading denotes parameter group (blue: oscillator, green: envelope, pink: filter, orange: misc). The oscillator group shows the largest absolute improvement; envelope parameters approach near-ceiling performance by V3.}
\label{fig:per_param}
\end{figure}

\noindent
The oscillator group carries the highest MAE in every iteration, with \texttt{osc2\_wave\_frame} and \texttt{osc1\_wave\_frame} the hardest parameters even in V3 ($0.1205$ and $0.1015$). The largest single-parameter improvement is \texttt{osc2\_transpose}, from $0.0902$ (V1) to $0.0622$ (V2) to $0.0308$ (V3). Envelope parameters record the lowest MAE throughout, with \texttt{env1\_decay} reaching $0.0103$ in V3. The discrete \texttt{osc1\_unison\_voices} parameter is the exception to the general downward trend, remaining at $0.1022$ in V3. These per-parameter patterns are analysed in Section~\ref{sec:discussion}.
